% Resumo e palavras-chave do projeto.

\begin{resumo}
    Este trabalho analisa a viabilidade técnica de um servidor \textit{Network Attached Storage} (\textit{NAS}) doméstico de baixo custo baseado em um \textit{Single Board Computer} \textit{Radxa ROCK 4B+}, com sistema base em cartão \textit{microSD} e dois \textit{Solid State Drives} (\textit{SSDs}) \textit{SATA} de 2,5 polegadas conectados por adaptadores \textit{USB-SATA}. A solução emprega \textit{Armbian}/\textit{Debian} como sistema operacional, \textit{OpenMediaVault} como camada de administração, \textit{SMB} para compartilhamento local, \textit{Docker} para a execução do \textit{Nextcloud}, do \textit{MariaDB} e do \textit{Redis}, e \textit{Tailscale} instalado diretamente no \textit{host} para acesso remoto privado. O problema de pesquisa consiste em verificar se essa arquitetura atende cenários domésticos de armazenamento, compartilhamento, sincronização e acesso remoto com desempenho, estabilidade e administrabilidade adequados. A avaliação combinou uma fase preliminar em rede cabeada \textit{Gigabit}, na qual o compartilhamento \textit{SMB} atingiu cerca de 76 \textit{MB/s} de escrita e 81 \textit{MB/s} de leitura, e uma fase consolidada com cenários por \textit{Wi-Fi}, com o \textit{NAS} em \textit{Ethernet} e o cliente em \textit{Wi-Fi}, e com acesso remoto por \textit{Tailscale} em rede móvel. No cenário híbrido, o \textit{SMB} alcançou cerca de 25 \textit{MB/s} de leitura e 18--19 \textit{MB/s} de escrita, e o \textit{Nextcloud}, 26--27 \textit{MB/s} de \textit{download} e 13--15 \textit{MB/s} de \textit{upload}; no acesso remoto, a leitura \textit{SMB} ficou próxima de 26 \textit{MB/s}. Os valores de referência de leitura, escrita e acesso remoto foram atendidos com o cliente cabeado, e o desempenho em \textit{Wi-Fi} foi limitado pela rede sem fio do cliente. O custo-base de R\$ 625,48, sem os \textit{SSDs}, foi cerca de 73\% inferior ao de \textit{NAS} comerciais de duas baias. Conclui-se que a arquitetura é tecnicamente viável para o contexto doméstico, com limitações ligadas à rede do cliente, ao barramento \textit{USB} e à configuração cuidadosa do \textit{Tailscale} e dos volumes \textit{Docker}.

    \palavraschave{\textit{NAS} doméstico; \textit{Single Board Computer}; \textit{OpenMediaVault}; \textit{Nextcloud}; \textit{Tailscale}}
\end{resumo}
